%% ============================================================================
%%  Lazy or Eager? When Post-Copy Process Restore Pays over a Network, and
%%  Streaming the Working Set.  IEEEtran conference format (venue TBD).
%%  Build: pdflatex regime && bibtex regime && pdflatex regime && pdflatex regime
%%  Data: eval/results/{regime,regime-stream,redis-v2,redis-same,pytorch-same,
%%        pi-stream,pi-stream-pfifo,pi-redis}/  (per-section pointers below)
%% ============================================================================
\documentclass[conference]{IEEEtran}

\usepackage[T1]{fontenc}
\usepackage[utf8]{inputenc}
\usepackage{microtype}
\usepackage{amsmath,amssymb}
\usepackage{booktabs}
\usepackage{multirow}
\usepackage{xcolor}
\usepackage{pgfplots}
\pgfplotsset{compat=1.18}
\usepackage{cite}
\usepackage[colorlinks=true,linkcolor=black,citecolor=black,urlcolor=blue!60!black]{hyperref}
\usepackage{xurl}
\usepackage{xspace}

% Categorical slots 1-4, same validated palette as the companion report
% (CVD-safe adjacent pairs; markers and tables carry identity as well).
\definecolor{cA}{HTML}{2A78D6}
\definecolor{cB}{HTML}{EB6834}
\definecolor{cC}{HTML}{1BAF7A}
\definecolor{cD}{HTML}{EDA100}

\newcommand{\uffd}{\texttt{userfaultfd}\xspace}
\newcommand{\us}{\,$\mu$s\xspace}
\newcommand{\mode}[1]{\textsf{#1}}

\title{Lazy or Eager? When Post-Copy Process Restore\\
Pays over a Network, and Streaming the Working Set}

\author{\IEEEauthorblockN{Utkarsh Maurya}
\IEEEauthorblockA{Independent Researcher\\
\href{mailto:utkarsh@kernex.sbs}{utkarsh@kernex.sbs}}}

\begin{document}
\emergencystretch=1em
\maketitle

\begin{abstract}
Restoring a checkpointed process on another machine can move the whole
image first (eager) or start at once and fetch pages on fault (lazy,
post-copy). Lazy restore is usually evaluated against an eager baseline that
reads its image from local disk, which hides the cost the network adds to
eager. We measure both honestly for CRIU process restore, with time to first
response (TTFR) as the metric, over a grid of round-trip times (0--2\,ms)
and bandwidths from 100\,Mb/s to 10\,Gb/s, for a PyTorch inference server and a
Redis-compatible key-value store. No fetch policy wins everywhere: demand
paging wins only when bandwidth is scarce and RTT is near zero, sequential
prefetch wins in the middle, and moving the whole image wins once
bandwidth${}\times{}$RTT is large, which for Redis is 10 of 12 grid points.
Each lazy policy pays roughly one RTT per demand fault, and that term
decides the map. We then remove it: the page server pushes a recorded
first-touch trace over its own connection at line rate while demand faults
keep a separate connection. With a trace recorded on an earlier run this
wins all 12 grid points for PyTorch (1.19--3.20$\times$ over the best
existing policy) but only 6 of 12 for Redis, whose heap layout is not
reproducible across runs on our x86 host. With a trace recorded on the
same image, as REAP does locally, it wins all 24 grid points
(1.22--5.75$\times$) and its TTFR comes within 1--10\,\% of the working
set's transfer time at 100\,Mb/s and 1\,Gb/s. On two Raspberry Pi~5s over
a real 1\,Gb/s link, streaming cuts TTFR by 1.5--3.0$\times$ for PyTorch
and 3.5--4.0$\times$ for Redis.
\end{abstract}

% ----------------------------------------------------------------------------
\section{Introduction}
% ----------------------------------------------------------------------------

Checkpoint/restore moves running processes between machines: for load
balancing, for draining a host, and increasingly for starting serverless
functions and agent sandboxes from a warm snapshot. CRIU~\cite{criu} offers
two ways to restore on a remote host. Eager restore copies the image and then
maps it. Lazy restore~\cite{criu-lazy} maps the process at once with its
memory registered with \uffd~\cite{userfaultfd}; a handler fetches each page
from a page server on the source the first time the process touches it.

The usual argument for lazy restore is that a process touches a small part of
its memory before it can do useful work. That argument compares lazy restore
over the network with eager restore from local disk. A real migration has no
local image: eager restore must first move it, at
$\text{image}/\text{bandwidth}$. Lazy restore pays one round trip per demand
fault instead. Which cost is smaller depends on the image size, the working
set, the bandwidth and the RTT, and no published study maps it for process
restore. REAP~\cite{ustiugov2021reap}, FaaSnap~\cite{ao2022faasnap} and
Spice~\cite{holmes2025spice} restore from local storage; AWS records chunk
access order for SnapStart and container loading~\cite{snapstart,
brooker2023ondemand} but has not published a regime analysis.

This paper makes four contributions.

\begin{enumerate}
\item \textbf{A regime map} (\S\ref{sec:map}). We measure eager-over-network,
demand paging, sequential/stride prefetch and an adaptive prefetcher across
four RTTs and three bandwidths. No policy wins everywhere. The losing
policies are not slow because of bandwidth: each demand fault costs about one
RTT, and the thousands of faults left after prefetching decide the outcome.
\item \textbf{Trace-ordered streaming} (\S\ref{sec:design}). The page server
pushes a recorded first-touch trace as one request over its own TCP
connection, at whatever rate TCP achieves; the handler installs pages as
they land. Demand faults keep their own connection. TTFR then depends on
the working set's transfer time plus one RTT per page the trace missed.
\item \textbf{Where traces come from matters} (\S\ref{sec:stream},
\S\ref{sec:misses}). A trace from an earlier run of the same program is
enough for PyTorch (98--99\,\% page overlap) but not for Redis on our x86
host (64.5\,\%), where allocator placement shifts whole pages between runs
even with ASLR off. A trace recorded on the same image removes misses
almost entirely and wins every grid point for both workloads.
\item \textbf{Real hardware} (\S\ref{sec:pi}). On two Raspberry Pi~5s
over a 1\,Gb/s link the results hold and are, if anything, stronger, and
one loopback effect, a large penalty for streaming the rest of the image
through a shared FIFO, turns out to be an emulation artifact.
\end{enumerate}

All measurements use \texttt{TCP\_NODELAY} on every page-server socket. A
companion report~\cite{pitfall} shows that without it, Nagle's algorithm
interacting with delayed ACKs inverts the measured effect of prefetching on
this system; we take the corrected setup as given here.

% ----------------------------------------------------------------------------
\section{Background and Related Work}
% ----------------------------------------------------------------------------

\textbf{Post-copy restore.} Post-copy migration starts the target at once and
pulls memory on demand; Hines et al.~\cite{hines2009postcopy} added
adaptive pre-paging for virtual machines on a gigabit LAN. CRIU implements
post-copy for processes: \texttt{criu restore --lazy-pages} registers the
restored memory with \uffd and hands the descriptor to a lazy-pages daemon,
which resolves each fault with \texttt{UFFDIO\_COPY}. An open CRIU pull
request streams images for post-copy~\cite{criu-stream-pr}.

\textbf{Snapshot cold starts.} REAP~\cite{ustiugov2021reap} observed that
serverless functions touch a stable working set after restore, records it on
the first restore, and prefetches it on later ones. FaaSnap~\cite{ao2022faasnap}
refines this for VM snapshots. Spice~\cite{holmes2025spice} shows that with
the right kernel primitives local process restore approaches warm-start
latency. These systems restore from local storage, where a page costs
microseconds, not an RTT. Firecracker~\cite{agache2020firecracker} documents
the \uffd handling rules a snapshot handler must follow~\cite{firecracker-uffd}.

\textbf{Network restore in industry.} AWS Lambda loads container images and
SnapStart snapshots on demand from a tiered cache and prefetches chunks in a
recorded access order~\cite{brooker2023ondemand, snapstart}. That is the
closest design to ours; the analysis of when it pays has not been published.

\textbf{Far-memory prefetching.} Leap~\cite{maruf2020leap} and
3PO~\cite{branner2022threepo} prefetch remote memory in steady state. Restore
is a different problem: every page starts remote, and the access stream before
the first response is short and mostly cold.

% ----------------------------------------------------------------------------
\section{System and Methodology}
\label{sec:setup}
% ----------------------------------------------------------------------------

\textbf{Pipeline.} The workload runs on the source, is checkpointed with
CRIU, and is restored on the target. The page server is a Python TCP server
that indexes the CRIU pagemap and serves 4\,KiB pages by virtual address. The
lazy-pages handler is a C daemon that owns the \uffd. It fetches the faulting
page synchronously and can issue prefetches asynchronously off the fault path.

\textbf{Policies.} We compare:
\begin{itemize}
\item \mode{full}: eager restore from images already on the target. This is a
floor, not a contender: it moves no data.
\item \mode{full-net}: eager restore after streaming the whole image from the
source over the same shaped link.
\item \mode{lazy}: demand paging only.
\item \mode{prefetch}: on each fault, also fetch the next pages
sequentially, or along a detected stride.
\item \mode{adaptive}: \mode{prefetch} with a controller that backs off when
prefetched pages go unused.
\item \mode{stream}, \mode{stream-pf}, \mode{stream-all}, \mode{stream-same}:
trace-ordered streaming (\S\ref{sec:design}).
\end{itemize}

\textbf{Workloads.} \emph{PyTorch}: a ResNet-18 inference server; the probe
writes a trigger file and waits for one inference result. Its image is
80.4k pages (329\,MB) on x86, of which 15.5k (63\,MB, 19\,\%) are touched
before the first result. \emph{Redis}: a key-value server holding 200k keys
of 1\,KiB; the probe issues a fixed set of 20k \texttt{MGET}s in batches of
100. On the x86 host the server binary is Valkey~9.1.2 (the distribution
installs it as \texttt{redis-server}); on the Pis it is Redis~8.0.2. The
x86 Redis trace is 29.8k pages (122\,MB). Probes reach Redis over a Unix
socket, so probe traffic is not shaped; the Pi kernel lacks
\texttt{unix\_diag}, so there the probe uses TCP over loopback.

\textbf{Metric.} TTFR is the time from the start of \texttt{criu restore}
until the probe's first response. For \mode{full-net} it includes the image
transfer.

\textbf{Testbeds.} \emph{x86}: one AMD Ryzen~7 7735HS laptop, Linux 7.2.8,
CRIU 4.2.1. Source and target share the host; the network is loopback
shaped with \texttt{netem}~\cite{hemminger2005netem}: delay $d$ on
\texttt{lo} (so RTT${}=2d$) and a rate cap. The grid is RTT
$\in\{0, 0.2, 1, 2\}$\,ms $\times$ rate $\in\{100\,\text{Mb/s},
1\,\text{Gb/s}, 10\,\text{Gb/s}\}$. \emph{Pi}: two Raspberry Pi~5s on a
direct 1\,GbE link (base RTT 164\us). Pi~A runs and restores the workload;
Pi~B runs the page server. \texttt{netem} delay goes on A's egress only,
so B's data keeps B's real qdisc. Pi~A runs Linux 6.18.59; PyTorch runs
used CRIU 4.1.1, Redis runs CRIU 4.2.1 (4.1.1 cannot dump the Redis TCP
socket on this kernel).

\textbf{Statistics.} Five iterations per cell, each with a fresh workload
and checkpoint. We report medians. For the existing policies the population
standard deviation is at most 3.6\,\% of the median (one outlier
iteration puts \mode{full}, a 200\,ms floor, at 10\,\%). Streaming with a
cross-run trace varies far more, up to 38\,\% for \mode{stream-all} and
31\,\% for Redis \mode{stream}, because the number of misses changes from
run to run (\S\ref{sec:misses}); every win we report for those modes holds
for all five iterations, not just the median. \mode{stream-same} stays
under 5\,\% except in six cells at 1--10\,Gb/s with one slow iteration.
Demand-only \mode{lazy} is skipped at 1 and 2\,ms RTT on x86, where it
exceeds the 30\,s probe timeout.

% ----------------------------------------------------------------------------
\section{The Regime Map}
\label{sec:map}
% ----------------------------------------------------------------------------

Table~\ref{tab:map} shows TTFR for the existing policies.
\emph{Data: \texttt{eval/results/regime/}, \texttt{redis-v2/}.}

\begin{table}[t]
\centering
\caption{TTFR (ms, median of 5) for existing policies on x86. Bold: best
contender. \mode{full} reads local images and is shown as a floor.}
\label{tab:map}
\setlength{\tabcolsep}{3.2pt}
\begin{tabular}{@{}rlrrrrr@{}}
\toprule
RTT & rate & full & full-net & lazy & prefetch & adaptive \\
\midrule
\multicolumn{7}{@{}l}{\emph{PyTorch} (image 329\,MB, working set 63\,MB)} \\
0   & 100M & 207 & 26664 & \textbf{6278} & 7706 & 7898 \\
0   & 1G   & 206 & 2899  & 1388 & \textbf{873} & 889 \\
0   & 10G  & 212 & 509   & 800  & \textbf{335} & 341 \\
0.2 & 100M & 210 & 26749 & 9405 & \textbf{7855} & 8060 \\
0.2 & 1G   & 213 & 2904  & 4544 & 1153 & \textbf{1134} \\
0.2 & 10G  & 222 & \textbf{521} & 3966 & 628 & 606 \\
1   & 100M & 222 & 26763 & --   & \textbf{9005} & 9096 \\
1   & 1G   & 208 & 2928  & --   & 2550 & \textbf{2286} \\
1   & 10G  & 212 & \textbf{549} & -- & 2188 & 1711 \\
2   & 100M & 220 & 26765 & --   & 10853 & \textbf{10464} \\
2   & 1G   & 210 & \textbf{2982} & -- & 4414 & 3726 \\
2   & 10G  & 215 & \textbf{627} & -- & 4231 & 3223 \\
\midrule
\multicolumn{7}{@{}l}{\emph{Redis} (Valkey 9.1.2, trace 122\,MB)} \\
0   & 100M & 179 & 22311 & \textbf{12331} & 23202 & 23906 \\
0   & 1G   & 180 & \textbf{2427} & 2687 & 2527 & 2515 \\
0   & 10G  & 190 & \textbf{440} & 1550 & 744 & 756 \\
0.2 & 100M & 182 & 22317 & \textbf{18192} & 24048 & 24656 \\
0.2 & 1G   & 177 & \textbf{2435} & 8822 & 3520 & 3480 \\
0.2 & 10G  & 185 & \textbf{436} & 7803 & 1909 & 1825 \\
1   & 100M & 175 & \textbf{22330} & -- & 28254 & 27257 \\
1   & 1G   & 175 & \textbf{2464} & -- & 8756 & 7775 \\
1   & 10G  & 174 & \textbf{478} & -- & 7594 & 5779 \\
2   & 100M & 177 & \textbf{22354} & -- & 34679 & 31579 \\
2   & 1G   & 172 & \textbf{2496} & -- & 15704 & 12980 \\
2   & 10G  & 182 & \textbf{533} & -- & 14800 & 10657 \\
\bottomrule
\end{tabular}
\end{table}

\textbf{No single policy wins.} For PyTorch there are three zones. Demand
paging wins only at 100\,Mb/s with zero RTT, where every extra byte is
expensive and faults are cheap. Prefetch or the adaptive variant wins the
middle of the grid. Eager \mode{full-net} wins once bandwidth${}\times{}$RTT
is large: at 10\,Gb/s from 0.2\,ms and at 1\,Gb/s from 2\,ms. For Redis
the eager zone covers 10 of 12 cells; demand paging wins the other two.

\textbf{Prefetch overfetches and misses the random workload.} For PyTorch,
prefetch serves about 22k pages against a 15.5k-page working set (42\,\%
extra), which is why it loses to demand paging at 100\,Mb/s. For Redis,
whose probe reads hash-table buckets and values at random, sequential
prefetch is worse than nothing at low RTT (23.2\,s vs 12.3\,s at
0/100\,Mb/s), and it still leaves 6--9k demand faults on the Pis
(\S\ref{sec:pi}).

\textbf{Each fault costs one RTT.} PyTorch under demand paging takes 15.5k
faults; raising the RTT from 0 to 0.2\,ms adds 3.2\,s, which is 15.5k
times 0.2\,ms. Prefetch leaves about 2k faults and gains 3.5\,s from 0 to
2\,ms RTT. Eager pays one bulk transfer and no per-fault RTT, which is
why it wins at high RTT despite moving five times more data than the
working set.

\textbf{The gap.} The cost a restore cannot avoid is moving the working set:
63\,MB is 5.1\,s at 100\,Mb/s, 0.51\,s at 1\,Gb/s and 51\,ms at 10\,Gb/s.
The best existing policy is 1.24$\times$ above that at 100\,Mb/s and
zero RTT, and 5.9$\times$ above it at 1\,Gb/s and 2\,ms. A policy that moved exactly the working set,
at line rate, without per-fault round trips, would close the gap.

% ----------------------------------------------------------------------------
\section{Trace-Ordered Streaming}
\label{sec:design}
% ----------------------------------------------------------------------------

\textbf{Recording.} A record run is a demand-only lazy restore in which the
handler appends each fault address to a trace file, in fault order. The
trace is the working set before the first response, ordered by first touch.
For traces reused across runs the workload is launched with ASLR off
(\texttt{setarch -R}) so that addresses are comparable between runs.

\textbf{Streaming.} At restore the handler opens a second connection to the
page server and sends one request: a magic marker, a count and the list of
trace addresses. The server answers with the pages back to back, in trace
order, without further requests. TCP alone sets the rate. A receive thread
installs each page with \texttt{UFFDIO\_COPY} as soon as its 4\,KiB are
complete and publishes how far into the trace it has got.

\textbf{Demand path.} Faults are served on the original connection, so a
fault never waits behind the stream at the client. The handler classifies
each fault against the stream: \emph{landed} (the stream already delivered
it; the fault raced the install), \emph{ahead} (in the trace, not yet
delivered) or \emph{miss} (not in the trace). Ahead and miss pages are
fetched on demand; an install that loses to the stream fails with
\texttt{EEXIST} and is ignored.

\textbf{Variants.} \mode{stream} uses a trace recorded on an earlier run.
\mode{stream-pf} adds sequential/stride prefetch, but only around misses.
\mode{stream-all} follows the trace with every other image page, ordered by
distance to the nearest traced page. \mode{stream-same} first runs an
untimed record pass on the same image and then restores it again, streaming
that trace; it models a system that records once per snapshot and restores
it many times, as REAP does.

\textbf{What it took.} A first version reused the prefetch thread: one 8\,B
request per page, up to 1024 in flight. The Python server then capped at
44k pages/s, and at 2\,ms RTT the request window let the process overtake
the stream, so most faults were ``ahead''. Server push removed both: the
server reaches about 275k pages/s on loopback at 10\,Gb/s and line rate at
1\,Gb/s, and the 0\,ms/10\,Gb/s PyTorch cell fell from 405 to 129\,ms.

% ----------------------------------------------------------------------------
\section{Streaming Results on x86}
\label{sec:stream}
% ----------------------------------------------------------------------------

\begin{table}[t]
\centering
\caption{Streaming vs the best existing policy on x86 (TTFR, ms). ``Best
old'': best of Table~\ref{tab:map}. ``Cross-run'': best of
\mode{stream}, \mode{stream-pf}, \mode{stream-all} with a trace from an
earlier run. Speedups of \mode{stream-same} over best old.}
\label{tab:stream}
\setlength{\tabcolsep}{3.0pt}
\begin{tabular}{@{}rlrrrr|rrr@{}}
\toprule
 & & \multicolumn{4}{c|}{PyTorch} & \multicolumn{3}{c}{Redis} \\
RTT & rate & old & cross & same & $\times$ & old & cross & same \\
\midrule
0   & 100M & 6278 & 5274 & \textbf{5129} & 1.22 & 12331 & 13234 & \textbf{9855} \\
0   & 1G   & 873  & 585  & \textbf{546}  & 1.60 & 2427 & 1814 & \textbf{1011} \\
0   & 10G  & 335  & 129  & \textbf{93}   & 3.60 & 440  & 331  & \textbf{131} \\
0.2 & 100M & 7855 & 5292 & \textbf{5136} & 1.53 & 18192 & 14632 & \textbf{9860} \\
0.2 & 1G   & 1134 & 682  & \textbf{544}  & 2.08 & 2435 & 2506 & \textbf{1024} \\
0.2 & 10G  & 521  & 163  & \textbf{92}   & 5.66 & 436  & 331  & \textbf{129} \\
1   & 100M & 9005 & 5751 & \textbf{5130} & 1.76 & 22330 & 22023 & \textbf{9875} \\
1   & 1G   & 2286 & 820  & \textbf{554}  & 4.13 & 2464 & 2569 & \textbf{1046} \\
1   & 10G  & 549  & 249  & \textbf{99}   & 5.55 & 478  & 357  & \textbf{140} \\
2   & 100M & 10464 & 5679 & \textbf{5147} & 2.03 & 22354 & 24372 & \textbf{9899} \\
2   & 1G   & 2982 & 1028 & \textbf{561}  & 5.32 & 2496 & 2647 & \textbf{1073} \\
2   & 10G  & 627  & 338  & \textbf{109}  & 5.75 & 533  & 437  & \textbf{214} \\
\bottomrule
\end{tabular}
\end{table}

\begin{figure*}[t]
\centering
\pgfplotsset{
  rmap/.style={
    width=0.47\textwidth, height=5.0cm,
    xmin=-0.1, xmax=2.1, log ticks with fixed point,
    xtick={0,0.2,1,2}, xlabel={RTT (ms)},
    ylabel={TTFR (ms), log scale},
    grid=major, grid style={gray!20},
    axis line style={gray!60}, tick style={gray!60},
    legend style={draw=none, fill=none, font=\footnotesize},
    legend cell align=left, every axis plot/.append style={line width=1.2pt, mark size=2.2pt},
  }}
\begin{tikzpicture}
\begin{axis}[rmap, ymode=log, title={PyTorch, 1\,Gb/s}, ymin=400, ymax=5000,
             ytick={500,1000,2000,4000},
             legend to name=rmaplegend, legend columns=5]
\addplot[color=cA, mark=square*] coordinates {(0,2899) (0.2,2904) (1,2928) (2,2982)};
\addlegendentry{eager over network (\mode{full-net})}
\addplot[color=cB, mark=triangle*] coordinates {(0,889) (0.2,1134) (1,2286) (2,3726)};
\addlegendentry{\mode{adaptive} prefetch}
\addplot[color=cC, mark=*] coordinates {(0,585) (0.2,682) (1,820) (2,1028)};
\addlegendentry{cross-run stream (best)}
\addplot[color=cD, mark=diamond*] coordinates {(0,546) (0.2,544) (1,554) (2,561)};
\addlegendentry{\mode{stream-same}}
\addplot[gray, dashed, line width=0.8pt, mark=none] coordinates {(-0.1,508) (2.1,508)};
\addlegendentry{working set / bandwidth}
\end{axis}
\end{tikzpicture}\hfill
\begin{tikzpicture}
\begin{axis}[rmap, ymode=log, title={Redis, 1\,Gb/s}, ymin=400, ymax=20000,
             ytick={500,1000,2000,5000,10000}]
\addplot[color=cA, mark=square*] coordinates {(0,2427) (0.2,2435) (1,2464) (2,2496)};
\addplot[color=cB, mark=triangle*] coordinates {(0,2515) (0.2,3480) (1,7775) (2,12980)};
\addplot[color=cC, mark=*] coordinates {(0,1814) (0.2,2506) (1,2569) (2,2647)};
\addplot[color=cD, mark=diamond*] coordinates {(0,1011) (0.2,1024) (1,1046) (2,1073)};
\addplot[gray, dashed, line width=0.8pt, mark=none] coordinates {(-0.1,976) (2.1,976)};
\end{axis}
\end{tikzpicture}
\\[2pt]\ref{rmaplegend}
\caption{TTFR against RTT at 1\,Gb/s on x86. Lazy prefetching grows by
about one RTT per remaining fault; eager pays for the whole image once;
same-image streaming stays within 10\,\% of the working set's transfer
time at every RTT.}
\label{fig:rtt}
\end{figure*}

Table~\ref{tab:stream} and Fig.~\ref{fig:rtt} summarize the streaming
modes. \emph{Data: \texttt{regime-stream/}, \texttt{redis-v2/},
\texttt{redis-same/}, \texttt{pytorch-same/}.}

\textbf{PyTorch, cross-run trace.} A stream mode wins all 12 cells, by
1.19--3.20$\times$ over the best existing policy. Plain \mode{stream} is
best at 100\,Mb/s, \mode{stream-pf} in the middle, and \mode{stream-all}
at 10\,Gb/s with RTT${}\ge{}$1\,ms. TTFR still rises with RTT (585 to
1028\,ms at 1\,Gb/s) because a few hundred faults miss the trace.

\textbf{Redis, cross-run trace.} Only 64.5\,\% of one recording's pages
appear in the other. A stream mode beats the best existing policy by
1.22--1.34$\times$ in 6 cells, ties within 2\,\% in one, and loses the
other five by 3--9\,\%: thousands of misses at one RTT each eat the gain.

\textbf{Same-image trace.} \mode{stream-same} wins all 24 cells: by
1.22--5.75$\times$ over the best existing policy for PyTorch and
1.25--3.41$\times$ for Redis, and by 1.03--3.10$\times$ and
1.34--2.57$\times$ over the best cross-run stream. Demand faults fall to
0--4 misses for PyTorch and under 120 faults of any class for Redis.
TTFR becomes nearly flat in RTT (PyTorch at 1\,Gb/s: 546 to 561\,ms over
0--2\,ms) and tracks bandwidth. It lands within 1--2\,\% of
$\text{working set}/\text{bandwidth}$ at 100\,Mb/s, 4--10\,\% at
1\,Gb/s, and 30--120\,ms above it at 10\,Gb/s, where restore overheads
dominate a 50--100\,ms transfer.

\textbf{The cost moved to the record pass.} The record pass is a
demand-only lazy restore. For PyTorch it takes 0.8--6.3\,s at zero RTT
and 32--38\,s at 2\,ms; for Redis about 65\,s at 2\,ms. \mode{stream-same}
is therefore a number for the second and later restores of a snapshot, not
the first. Systems that restore one snapshot many times, such as
serverless platforms and agent sandboxes, amortize it; a one-off migration
cannot, and must use a cross-run trace or an existing policy.

% ----------------------------------------------------------------------------
\section{Why Cross-Run Traces Miss}
\label{sec:misses}
% ----------------------------------------------------------------------------

Every miss costs one RTT on the demand connection, so the trace's coverage
of the next run decides how close a cross-run stream gets to
\mode{stream-same}.

\textbf{PyTorch: fresh heap.} Two PyTorch recordings share 98.0--99.0\,\%
of pages with a median rank shift of 4. The differences sit in the
\texttt{brk} heap: one run had a 617-page heap region the other lacked.
Even the union of three runs' traces leaves 85--690 misses against a
fourth. Every run allocates some fresh heap that no earlier trace can name.

\textbf{Redis on x86: layout jitter.} The two x86 Redis traces have the
same number of pages in each 256\,MB region (27,210 vs 27,212), but the long
runs of pages sit a few pages apart: a 1313-page run starts at
\texttt{\ldots2f8e000} in one trace and \texttt{\ldots2f7b000} in the other.
The process maps are identical under \texttt{setarch -R}; placement inside
the allocator's arenas differs. We started the server three times with
ASLR off, loaded the same 200k keys, and read the addresses of 201 sampled
keys with \texttt{DEBUG OBJECT}:

\begin{center}
\small
\begin{tabular}{@{}lr@{}}
\toprule
setup & keys at the same address \\
\midrule
x86, Valkey 9.1.2, jemalloc 5.4.0 & 64--183 / 201 \\
\quad hash seed fixed & 154--182 / 201 \\
\quad jemalloc 5.3.1 swapped in & 142--156 / 201 \\
\quad pinned to 4 CPUs / 1 CPU & 65--135 / 25--146 \\
Pi, Valkey 8.1.1, jemalloc 5.3.0 & 197--201 / 201 \\
Pi, Redis 8.0.2 & 193--199 / 201 \\
\bottomrule
\end{tabular}
\end{center}

\noindent The mismatches are whole-page offsets (4, 8 or 20\,KB) that
appear and vanish along the insertion order rather than accumulating,
consistent with freed multi-page extents being reused in a timing-dependent
order. Valkey's random hash seed contributes (its hash table allocates
overflow buckets depending on where keys hash), and the allocator build
contributes; disabling active rehashing or jemalloc's decay-based purging
does not help. We did not isolate the rest, nor separate Valkey 9.1 from
8.1. The same server family is reproducible on the Pi, where the
cross-run trace overlaps 98.8\,\% and cross-run streaming comes within
4\,\% of \mode{stream-same} (\S\ref{sec:pi}).

\textbf{Implication.} Cross-run traces work only for programs whose heap
layout is reproducible, and that is a property of the build and platform,
not just of the program. A same-image trace is immune: its addresses come
from the image itself.

% ----------------------------------------------------------------------------
\section{Streaming the Rest of the Image}
\label{sec:rest}
% ----------------------------------------------------------------------------

\mode{stream-all} covers misses by streaming every remaining page after the
trace. It cuts PyTorch faults to about 50 and wins at 10\,Gb/s with high
RTT, but on loopback it loses badly at 100\,Mb/s: 55 misses cost 12.4\,s,
about 225\,ms each, because each demand reply queues behind the stream at
the shaped bottleneck. Delay-based congestion control (Vegas) on the stream
socket or pacing the rest at 0.7--0.85$\times$ the measured delivery rate
roughly halves the damage; TCP-LP does nothing, because demand traffic is
too small to count as a competing flow. No endpoint-only variant matches
plain \mode{stream} at 100\,Mb/s.

On the Pis, the same comparison shows a penalty of about 1.1\,ms per miss,
not 225\,ms, and it is unchanged when Pi~B's root qdisc is switched from
\texttt{fq\_codel} to a plain \texttt{pfifo} (Table~\ref{tab:pfifo}). We
believe TCP Small Queues bounds how much of the sender's data waits in the
qdisc and driver, so the queue a demand reply waits behind is a millisecond
or two at line rate, whereas \texttt{netem} on loopback holds a deep emulated
queue. We did not verify the TSQ explanation directly. The loopback FIFO
penalty should be read as an emulation artifact, not an estimate for a real
link.

\begin{table}[t]
\centering
\caption{Pi PyTorch, \mode{stream} and \mode{stream-all} with Pi~B's qdisc
\texttt{fq\_codel} vs \texttt{pfifo}: TTFR (ms) and median ms per miss.}
\label{tab:pfifo}
\setlength{\tabcolsep}{4pt}
\begin{tabular}{@{}lrrrrr@{}}
\toprule
mode & delay & fq\_codel & pfifo & ms/miss fq & ms/miss pfifo \\
\midrule
stream     & 0    & 728  & 718  & 0.88 & 0.91 \\
stream     & 1000 & 1127 & 1071 & 1.58 & 1.62 \\
stream     & 2000 & 1594 & 1405 & 2.48 & 2.54 \\
stream-all & 0    & 1145 & 1127 & 1.99 & 2.25 \\
stream-all & 1000 & 1249 & 1178 & 2.50 & 2.63 \\
stream-all & 2000 & 1372 & 1292 & 3.53 & 3.56 \\
\bottomrule
\end{tabular}
\end{table}

% ----------------------------------------------------------------------------
\section{Real Hardware: Two Raspberry Pi 5s}
\label{sec:pi}
% ----------------------------------------------------------------------------

\begin{table}[t]
\centering
\caption{Pi TTFR (ms, median of 5) over a real 1\,GbE link. RTT is the
164\us base plus the added delay. Speedup: best streaming mode over the
best existing policy. PyTorch: image 74.7k pages, CRIU 4.1.1. Redis 8.0.2:
image 437\,MB, trace 111\,MB, CRIU 4.2.1.}
\label{tab:pi}
\setlength{\tabcolsep}{4pt}
\begin{tabular}{@{}lrrr@{}}
\toprule
mode & 0.16\,ms & 1.16\,ms & 2.16\,ms \\
\midrule
\multicolumn{4}{@{}l}{\emph{PyTorch}} \\
full (floor)   & 277  & 269  & 275 \\
full-net       & 2939 & 2937 & 2947 \\
prefetch       & 1023 & 2850 & 5014 \\
adaptive       & 1029 & 2569 & 4226 \\
stream         & 728  & 1127 & 1594 \\
stream-pf      & \textbf{700} & \textbf{867} & \textbf{1076} \\
stream-all     & 1145 & 1249 & 1372 \\
speedup        & 1.46$\times$ & 2.96$\times$ & 2.74$\times$ \\
\midrule
\multicolumn{4}{@{}l}{\emph{Redis}} \\
full (floor)   & 259  & 260  & 258 \\
full-net       & 4059 & 4064 & 4073 \\
lazy           & 4586 & --   & -- \\
prefetch       & 3606 & 9989 & 17709 \\
adaptive       & 3616 & 8218 & 13433 \\
stream         & 1018 & 1155 & 1359 \\
stream-pf      & 1021 & 1049 & 1058 \\
stream-all     & 1317 & 1715 & 1742 \\
stream-same    & \textbf{1002} & \textbf{1006} & \textbf{1014} \\
speedup        & 3.60$\times$ & 4.04$\times$ & 4.02$\times$ \\
\bottomrule
\end{tabular}
\end{table}

Table~\ref{tab:pi} repeats the comparison on two Raspberry Pi~5s.
\emph{Data: \texttt{pi-stream/}, \texttt{pi-stream-pfifo/},
\texttt{pi-redis/}.}

\textbf{The loopback result holds.} For PyTorch, \mode{stream-pf} wins
every delay by 1.46--2.96$\times$, close to the x86 1\,Gb/s cells
(1.49--2.90$\times$), and the existing policies order the same way:
prefetch wins at low RTT and eager from about 2\,ms. Each miss costs about
one RTT plus 0.7\,ms of service time on the Pi; prefetching around misses
removes a third of them, which is where \mode{stream-pf}'s lead over
\mode{stream} comes from.

\textbf{Redis gains most.} Redis touches 111\,MB of a 437\,MB image before
its first reply. Eager moves all of it (about 3.5\,s at line rate);
streaming moves the working set. Streaming wins by 3.5--4.0$\times$ at
every delay. \mode{stream-same} sits at the bandwidth floor: 111\,MB at
1\,Gb/s is about 0.9\,s, it took 1.00--1.01\,s, the process overtook the
stream about 400 times (``ahead'' faults) and only 7 pages missed. Here the
cross-run trace overlaps 98.8\,\%, so \mode{stream-pf} comes within 2--4\,\%
of \mode{stream-same}; the gap that mattered on x86 is absent.

\textbf{Old lazy policies collapse with RTT.} Redis reads its hash table at
random, so prefetch guesses badly and still leaves 6--9k faults: the
adaptive prefetcher reaches 13.4\,s at 2.16\,ms RTT, more than three times
the eager time.

% ----------------------------------------------------------------------------
\section{Implementation Pitfalls}
% ----------------------------------------------------------------------------

Three problems cost us days and are easy to repeat.

\textbf{\texttt{UFFDIO\_COPY} returns \texttt{EAGAIN} and leaves the thread
asleep.} While a \texttt{REMOVE} or \texttt{UNMAP} event is pending (jemalloc
returns memory with \texttt{MADV\_DONTNEED}), the kernel refuses
\texttt{UFFDIO\_COPY} and \texttt{UFFDIO\_ZEROPAGE} with \texttt{EAGAIN}
and does not wake the faulting thread. A handler that treats the error as
``someone else will install it'' hangs the process; our Redis runs stalled
for hours. The fix is to issue \texttt{UFFDIO\_WAKE} on that range after
\texttt{EAGAIN} so the thread refaults.

\textbf{Image formats drift.} CRIU 4.2 writes the pagemap entry's page count
as a new 64-bit field and leaves the old one zero; an older \texttt{pycriu}
then reads every entry as empty and the page server serves zeros. We fall
back to decoding the field by hand.

\textbf{Restoring an image twice.} \mode{stream-same} restores one image
twice. The first restored process appends to its stdout log, and CRIU
refuses to restore a process whose open file has changed size; CRIU 4.2
also refuses to overwrite an existing pidfile. The harness truncates the
log to its dump-time size and removes the pidfile between passes.

% ----------------------------------------------------------------------------
\section{Discussion and Limitations}
% ----------------------------------------------------------------------------

\textbf{What to deploy.} If a snapshot is restored more than once, record
its first-touch trace on the first restore and stream it on later ones:
TTFR then approaches the working set's transfer time at every RTT we
tested. Cross-run traces are worth trying only where the program's heap
layout is reproducible; checking two recordings' overlap is cheap and
predicts the outcome. Without a trace, the regime map applies: demand
paging at low bandwidth and near-zero RTT, prefetching in the middle, and
eager transfer once bandwidth${}\times{}$RTT is large or the access
pattern is random.

\textbf{Limitations.}
\begin{itemize}
\item Two workloads. Both touch a quarter of their image or less before
the first response (19\,\% for PyTorch on x86, 25\,\% for Redis on the
Pi); a workload that touches most of its image
leaves streaming little to save over eager.
\item x86 results use \texttt{netem} on loopback, with source and target on
one host. The Pi results show one loopback effect (\S\ref{sec:rest}) does
not transfer; others might not either.
\item The page server is single-threaded Python. At 10\,Gb/s its throughput
and restore overheads, not the network, set the remaining 30--120\,ms.
\item Five iterations per cell. Cross-run streaming is noisy
(standard deviation up to 38\,\% of the median); more iterations would
narrow its medians, though not change which policy wins.
\item TTFR only. We do not measure throughput after the first response,
during which a lazy restore is still paging.
\item The Pi PyTorch and Redis runs used different CRIU versions (4.1.1 and
4.2.1); we did not run \mode{stream-same} for PyTorch on the Pi.
\item The cause of x86 heap-layout jitter is only partly identified.
\end{itemize}

% ----------------------------------------------------------------------------
\section{Conclusion}
% ----------------------------------------------------------------------------

Against an honest eager baseline that pays for moving its image, lazy
restore is not a free win: for a random-access key-value store, eager
transfer beats every lazy policy at 10 of 12 network settings we tested.
What decides the map is one round trip per demand fault. Streaming a
recorded first-touch trace by server push removes most of those round
trips. With a trace recorded on the same snapshot, it beats every existing
policy at all 24 settings, by up to 5.75$\times$, and stays within 10\,\%
of the time to move the working set at 100\,Mb/s and 1\,Gb/s. On real
hardware it cuts TTFR by up to 4$\times$.

\section*{Artifact Availability}
Code, harnesses and raw results are in the DistriProc repository
(\texttt{src/}, \texttt{eval/}); each section names its result directory
under \texttt{eval/results/}.

\bibliographystyle{IEEEtran}
\bibliography{references}

\end{document}
